%=========================================================
%               CLASSE DU DOCUMENT
%=========================================================
\documentclass[12pt,a4paper]{article}

%=========================================================
%         ENCODAGE ET POLICES DE CARACTÈRES
%=========================================================
\usepackage[utf8]{inputenc}    % Encodage UTF-8 (inutile avec XeLaTeX/LuaLaTeX)
\usepackage[T1]{fontenc}       % Meilleure gestion des caractères accentués
\usepackage{lmodern}           % Police Latin Modern
\usepackage{times}             % Police Times (à supprimer si tu préfères Latin Modern)

%=========================================================
%                MATHÉMATIQUES
%=========================================================
\usepackage{amsmath}           % Environnements mathématiques avancés
\usepackage{amssymb}           % Symboles mathématiques supplémentaires
\usepackage{amsfonts}          % Polices mathématiques
\usepackage{mathrsfs}          % Police calligraphique (\mathscr)
\usepackage{tkz-tab}           % Tableaux de signes et de variations
\everymath{\displaystyle}      % Affiche toutes les formules en style "display"

%=========================================================
%               MISE EN PAGE
%=========================================================
\usepackage[
    left=1cm,
    right=1cm,
    top=1cm,
    bottom=1.7cm
]{geometry}                    % Marges du document

\usepackage{multicol}          % Texte sur plusieurs colonnes
\usepackage{float}             % Placement précis des figures et tableaux

%=========================================================
%            TABLEAUX ET LISTES
%=========================================================
\usepackage{array}             % Colonnes de tableaux personnalisées
\usepackage{multirow}          % Fusion de lignes dans les tableaux
\usepackage{enumitem}          % Personnalisation des listes

\usepackage{tikz}
\usetikzlibrary{angles,quotes,arrows.meta}
%=========================================================
%            COULEURS ET BOÎTES
%=========================================================
\usepackage[most]{tcolorbox}   % Boîtes colorées très personnalisables
\tcbuselibrary{skins,hooks}    % Bibliothèques supplémentaires de tcolorbox

\usepackage{mdframed}          % Cadres autour de paragraphes
\usepackage{varwidth}          % Largeur automatique des boîtes

%=========================================================
%                DESSINS ET GRAPHIQUES
%=========================================================
\usepackage{tikz}              % Dessins vectoriels
\usetikzlibrary{calc}
\usetikzlibrary{
    arrows,
    shapes.geometric,
    fit,
    patterns
}

\usepackage{pgfplots}          % Tracés de fonctions et graphiques
\pgfplotsset{compat=1.15}

%=========================================================
%              LIENS HYPERTEXTES
%=========================================================
\usepackage[
    colorlinks=true,
    linkcolor=blue,
    urlcolor=blue,
    citecolor=blue
]{hyperref}

%=========================================================
%          ÉLÉMENTS EN ARRIÈRE-PLAN
%=========================================================
\usepackage{eso-pic}           % Ajout de filigranes ou d'éléments sur chaque page

%=========================================================
%              SYMBOLES DIVERS
%=========================================================
\usepackage{pifont}            % Symboles (✓, ✗, etc.)

\DeclareUnicodeCharacter{2460}{\text{\ding{172}}} % Pour ①
\DeclareUnicodeCharacter{2461}{\text{\ding{173}}} % Pour ②
%=========================================================
%      PERSONNALISATION DES LISTES NUMÉROTÉES
%=========================================================

\newcommand\circitem[1]{%
\tikz[baseline=(char.base)]{
\node[circle,draw=gray,fill=red!55,
minimum size=1.2em,inner sep=0] (char) {#1};}}

\newcommand\boxitem[1]{%
\tikz[baseline=(char.base)]{
\node[fill=cyan,
minimum size=1.2em,inner sep=0] (char) {#1};}}

\setlist[enumerate,1]{label=\protect\circitem{\arabic*}}
\setlist[enumerate,2]{label=\protect\boxitem{\alph*}}

%=========================================================
%            COMMANDES PERSONNALISÉES
%=========================================================

% Soulignement coloré
\newcommand{\myul}[2][black]{%
\setulcolor{#1}\ul{#2}\setulcolor{black}}

% Tabulation horizontale
\newcommand\tab[1][1cm]{\hspace*{#1}}

%=========================================================
%             COMPTEURS PERSONNALISÉS
%=========================================================

%---------- Exemple ----------
\newcounter{exemple}

\newcommand{\exemple}{%
\refstepcounter{exemple}
\textbf{\textcolor{orange}{Exemple \theexemple : }}
\ignorespaces}

%---------- Solution ----------
\newcounter{solution}

\newcommand{\solution}{%
\refstepcounter{solution}
\textbf{\textcolor{orange}{Solution \thesolution : }}
\ignorespaces}

%---------- Exercices ----------
\definecolor{myorange}{rgb}{1.0,0.8,0.0}

\newcounter{exerciceapp}

\newcommand{\exerciceapp}{%
\refstepcounter{exerciceapp}
\textbf{\textcolor{myorange}
{Exercice d'application \theexerciceapp : }}
\ignorespaces}

%---------- Corrections ----------
\newcounter{correction}

\newcommand{\correction}{%
\refstepcounter{correction}
\textbf{\textcolor{myorange}
{Correction \thecorrection : }}
\ignorespaces}

%---------- Remarques ----------
\definecolor{myorange1}{rgb}{1.0,1.0,0.0}
\usepackage{enumitem}

\definecolor{myred}{RGB}{180, 40, 40}
\definecolor{myblue}{RGB}{20, 50, 120}
\newcounter{remarque}

\newcommand{\remarque}{%
\refstepcounter{remarque}
\textbf{\textcolor{myorange1}
{Remarque \theremarque : }}
\ignorespaces}

%=========================================================
%                 FILIGRANE
%=========================================================

\AddToShipoutPicture{
    \AtTextCenter{
        \makebox[0pt]{
            \rotatebox{45}{
                \textcolor[gray]{0.6}{
                    \fontsize{5cm}{5cm}\selectfont PGB
                }
            }
        }
    }
}

\begin{document}
% En-tête personnalisée
\begin{center}
\Large\textbf{\underline{\textcolor{red}{CALCUL DANS $\mathbb{R}$}}}\\[-0.1cm]
    \normalsize\textbf{Prof : M. BA} \hfill \textbf{Classe : 2ndL}\\[-0.1cm]
    \textbf{Année scolaire : 2025 -- 2026}
\end{center}
%=========================================================

% --- I) CALCULS SUR LES QUOTIENTS ---
\section*{\color{myred} I) CALCULS SUR LES QUOTIENTS}

\subsection*{\color{myred} 1) Règles de calculs}

\subsubsection*{\color{myred} a) Quotient de deux réels}
Soient $a$ et $b$ deux réels, avec $b \neq 0$. Le quotient de $a$ par $b$, noté $\dfrac{a}{b}$, est le nombre réel qui, multiplié par $b$, donne $a$. Autrement dit :
\[
\dfrac{a}{b} = a \times \dfrac{1}{b}
\]

\subsubsection*{\color{myred} b) Propriétés}
Pour tous réels $a$, $b$, $c$ et $d$ (avec $b \neq 0$ et $d \neq 0$), on a les propriétés suivantes :
\begin{itemize}
    \item \textbf{Égalité des quotients :} $\dfrac{a}{b} = \dfrac{c}{d} \iff a \times d = b \times c$ (produit en croix).
    \item \textbf{Simplification :} $\dfrac{a \times c}{b \times c} = \dfrac{a}{b}$ pour $c \neq 0$.
    \item \textbf{Addition/Soustraction :} $\dfrac{a}{b} \pm \dfrac{c}{d} = \dfrac{a \times d \pm b \times c}{b \times d}$ (mettre au même dénominateur).
    \item \textbf{Multiplication :} $\dfrac{a}{b} \times \dfrac{c}{d} = \dfrac{a \times c}{b \times d}$.
    \item \textbf{Division :} $\dfrac{\dfrac{a}{b}}{\dfrac{c}{d}} = \dfrac{a}{b} \times \dfrac{d}{c} = \dfrac{a \times d}{b \times c}$ pour $c \neq 0$.
\end{itemize}

\subsubsection*{\color{myred} c) Application}
Simplifier au maximum l'expression suivante :
\[
A = \frac{3}{4} + \frac{5}{6} - \frac{7}{12}
\]
On met tout sur le même dénominateur (12 est le PPCM de 4, 6 et 12) :
\[
A = \frac{3 \times 3}{4 \times 3} + \frac{5 \times 2}{6 \times 2} - \frac{7}{12}
= \frac{9}{12} + \frac{10}{12} - \frac{7}{12}
= \frac{9+10-7}{12} = \frac{12}{12} = 1
\]

\subsection*{\color{myred} 2) IDENTITES REMARQUABLES}

\subsubsection*{\color{myred} a) Rappels}
Pour tous réels $a$ et $b$, on a les trois identités remarquables fondamentales :
\begin{enumerate}
    \item $(a+b)^2 = a^2 + 2ab + b^2$ (carré d'une somme).
    \item $(a-b)^2 = a^2 - 2ab + b^2$ (carré d'une différence).
    \item $(a-b)(a+b) = a^2 - b^2$ (produit de la somme par la différence).
\end{enumerate}

\begin{itemize}
    \item \textbf{Exemples :}
    \begin{itemize}
        \item Développer $(2x+3)^2 = (2x)^2 + 2 \times (2x) \times 3 + 3^2 = 4x^2 + 12x + 9$.
        \item Factoriser $x^2 - 16 = x^2 - 4^2 = (x-4)(x+4)$.
    \end{itemize}
\end{itemize}

\subsubsection*{\color{myred} b) Autres identités remarquables}
On généralise aux cubes et aux puissances $n$ :
\begin{itemize}
    \item $(a+b)^3 = a^3 + 3a^2b + 3ab^2 + b^3$.
    \item $(a-b)^3 = a^3 - 3a^2b + 3ab^2 - b^3$.
    \item $a^3 + b^3 = (a+b)(a^2 - ab + b^2)$.
    \item $a^3 - b^3 = (a-b)(a^2 + ab + b^2)$.
    \item Pour $n$ entier naturel non nul : $a^n - b^n = (a-b)(a^{n-1} + a^{n-2}b + \cdots + ab^{n-2} + b^{n-1})$.
\end{itemize}

\subsubsection*{\color{myred} c) Application}
Factoriser l'expression $B = 8x^3 - 27$.
On reconnaît une différence de cubes : $8x^3 = (2x)^3$ et $27 = 3^3$.
Donc
\[
B = (2x)^3 - 3^3 = (2x-3)\left((2x)^2 + (2x)\times 3 + 3^2\right) = (2x-3)(4x^2 + 6x + 9).
\]


% --- II) PUISSANCE D'UN REEL ---
\section*{\color{myred} II) PUISSANCE D'UN REEL}

\subsection*{\color{myred} 1) Définition}
Soit $a$ un réel et $n$ un entier naturel non nul. La puissance $n$-ième de $a$, notée $a^n$, est le produit de $n$ facteurs tous égaux à $a$ :
\[
a^n = \underbrace{a \times a \times \cdots \times a}_{n \text{ facteurs}}
\]
Par convention : $a^0 = 1$ (pour $a \neq 0$). Pour $a \neq 0$, on définit les puissances négatives : $a^{-n} = \dfrac{1}{a^n}$.

\begin{itemize}
    \item \textbf{Exemples :}
    \begin{itemize}
        \item $2^3 = 2 \times 2 \times 2 = 8$.
        \item $(-3)^2 = (-3) \times (-3) = 9$.
        \item $5^{-2} = \dfrac{1}{5^2} = \dfrac{1}{25}$.
        \item $10^0 = 1$.
    \end{itemize}
\end{itemize}

\subsection*{\color{myred} 2) Propriétés}
Pour tous réels $a$ et $b$ non nuls, et pour tous entiers relatifs $m$ et $n$, on a :
\begin{itemize}
    \item $a^m \times a^n = a^{m+n}$.
    \item $\dfrac{a^m}{a^n} = a^{m-n}$.
    \item $(a^m)^n = a^{m \times n}$.
    \item $(a \times b)^n = a^n \times b^n$.
    \item $\left(\dfrac{a}{b}\right)^n = \dfrac{a^n}{b^n}$.
\end{itemize}

\subsection*{\color{myred} 3) Application}
Simplifier l'expression suivante en une seule puissance :
\[
C = \frac{3^5 \times (3^{-2})^4}{3^{-3} \times 3^2}
\]
On applique les règles :
\[
C = \frac{3^{5} \times 3^{-8}}{3^{-1}} = \frac{3^{5-8}}{3^{-1}} = \frac{3^{-3}}{3^{-1}} = 3^{-3 - (-1)} = 3^{-2} = \frac{1}{9}.
\]


% --- III) CALCULS SUR LES RADICAUX ---
\section*{\color{myred} III) CALCULS SUR LES RADICAUX}

\subsection*{\color{myred} 1) Définition}
Soit $a$ un réel positif. La racine carrée de $a$, notée $\sqrt{a}$, est le nombre réel positif dont le carré est égal à $a$ :
\[
(\sqrt{a})^2 = a \quad \text{et} \quad \sqrt{a} \geq 0.
\]
De manière générale, pour $n$ entier naturel non nul, la racine $n$-ième de $a$ (avec $a \geq 0$) est le nombre positif dont la puissance $n$-ième vaut $a$.

\begin{itemize}
    \item \textbf{Exemples :}
    \begin{itemize}
        \item $\sqrt{16} = 4$ car $4^2 = 16$.
        \item $\sqrt[3]{8} = 2$ car $2^3 = 8$.
        \item $\sqrt{2} \approx 1.414$.
    \end{itemize}
\end{itemize}

\subsection*{\color{myred} 2) Propriétés}
Pour $a \geq 0$ et $b \geq 0$, on a :
\begin{itemize}
    \item $\sqrt{a \times b} = \sqrt{a} \times \sqrt{b}$.
    \item $\sqrt{\dfrac{a}{b}} = \dfrac{\sqrt{a}}{\sqrt{b}}$ pour $b > 0$.
    \item $\sqrt{a^2} = |a|$ (valeur absolue de $a$).
    \item Pour les racines $n$-ièmes : $\sqrt[n]{a^n} = a$ si $a \geq 0$, et pour $n$ impair, $\sqrt[n]{a^n} = a$ pour tout $a$ réel.
\end{itemize}
\begin{itemize}
    \item \textbf{Remarque :} On ne peut pas simplifier $\sqrt{a+b}$ en $\sqrt{a} + \sqrt{b}$ en général. Par exemple, $\sqrt{9+16} = \sqrt{25} = 5$, alors que $\sqrt{9} + \sqrt{16} = 3+4 = 7$.
\end{itemize}

\subsection*{\color{myred} 3) Application}
Simplifier l'expression $D = \sqrt{50} - 3\sqrt{18} + 2\sqrt{8}$.
On commence par décomposer les radicandes en produits de carrés parfaits :
\[
\sqrt{50} = \sqrt{25 \times 2} = 5\sqrt{2}, \quad
\sqrt{18} = \sqrt{9 \times 2} = 3\sqrt{2}, \quad
\sqrt{8} = \sqrt{4 \times 2} = 2\sqrt{2}.
\]
Donc
\[
D = 5\sqrt{2} - 3 \times 3\sqrt{2} + 2 \times 2\sqrt{2}
= 5\sqrt{2} - 9\sqrt{2} + 4\sqrt{2} = (5-9+4)\sqrt{2} = 0.
\]


% --- IV) INTERVALLES DANS R ---
\section*{\color{myred} IV) INTERVALLES DANS $\mathbb{R}$}

\subsection*{\color{myred} 1) Définition}
Un intervalle de $\mathbb{R}$ est un ensemble de nombres réels compris entre deux bornes (éventuellement infinies). On distingue les intervalles ouverts (bornes exclues), fermés (bornes incluses), ou semi-ouverts.

\subsection*{\color{myred} 2) Présentation des différents types d'intervalles}

\subsubsection*{\color{myred} a) Intervalles bornés}
Soient $a$ et $b$ deux réels avec $a < b$.
\begin{itemize}
    \item $[a, b] = \{ x \in \mathbb{R} \mid a \leq x \leq b \}$ (fermé).
    \item $]a, b[ = \{ x \in \mathbb{R} \mid a < x < b \}$ (ouvert).
    \item $[a, b[ = \{ x \in \mathbb{R} \mid a \leq x < b \}$ (fermé à gauche, ouvert à droite).
    \item $]a, b] = \{ x \in \mathbb{R} \mid a < x \leq b \}$ (ouvert à gauche, fermé à droite).
\end{itemize}

\begin{itemize}
    \item \textbf{Exemples :}
    \begin{itemize}
        \item $[2, 5]$ contient tous les nombres de 2 à 5 inclus.
        \item $]-3, 1]$ contient les nombres strictement plus grands que -3 et inférieurs ou égaux à 1.
    \end{itemize}
\end{itemize}

\subsubsection*{\color{myred} b) Intervalles non bornés}
\begin{itemize}
    \item $[a, +\infty[ = \{ x \in \mathbb{R} \mid x \geq a \}$.
    \item $]a, +\infty[ = \{ x \in \mathbb{R} \mid x > a \}$.
    \item $]-\infty, b] = \{ x \in \mathbb{R} \mid x \leq b \}$.
    \item $]-\infty, b[ = \{ x \in \mathbb{R} \mid x < b \}$.
    \item $]-\infty, +\infty[ = \mathbb{R}$.
\end{itemize}

\subsection*{\color{myred} 3) Intersection et réunion de deux intervalles}

\subsubsection*{\color{myred} a) Intersection de deux intervalles}
L'intersection de deux intervalles $I$ et $J$, notée $I \cap J$, est l'ensemble des réels qui appartiennent à la fois à $I$ et à $J$.

\begin{itemize}
    \item \textbf{Définition et Exemples :}
    \begin{itemize}
        \item $[1, 4] \cap [2, 6] = [2, 4]$.
        \item $]-\infty, 3] \cap ]0, +\infty[ = ]0, 3]$.
        \item $[1, 2] \cap [3, 4] = \emptyset$ (ensemble vide, car aucun réel commun).
    \end{itemize}
\end{itemize}

\subsubsection*{\color{myred} b) Réunion de deux intervalles}
La réunion de deux intervalles $I$ et $J$, notée $I \cup J$, est l'ensemble des réels qui appartiennent à $I$ ou à $J$ (ou aux deux).

\begin{itemize}
    \item \textbf{Définition et Exemples :}
    \begin{itemize}
        \item $[1, 4] \cup [2, 6] = [1, 6]$ (car ils se chevauchent).
        \item $]-\infty, 1] \cup [3, +\infty[ = ]-\infty, 1] \cup [3, +\infty[$ (pas d'écriture plus simple car il y a un "trou" entre 1 et 3).
        \item $[1, 2] \cup [3, 4] = [1, 2] \cup [3, 4]$ (pas d'intervalle unique).
    \end{itemize}
\end{itemize}


% --- V) VALEUR ABSOLUE ---
\section*{\color{myred} V) VALEUR ABSOLUE}

\subsection*{\color{myred} 1) Définition et Exemples}
La valeur absolue d'un réel $x$, notée $|x|$, est la distance de $x$ à 0 sur la droite réelle. Elle est toujours positive ou nulle.
\[
|x| =
\begin{cases}
x & \text{si } x \geq 0, \\
-x & \text{si } x < 0.
\end{cases}
\]

\begin{itemize}
    \item \textbf{Exemples :}
    \begin{itemize}
        \item $|5| = 5$ ; $|-3| = 3$ ; $|0| = 0$.
        \item $|x| = 2$ signifie $x = 2$ ou $x = -2$.
        \item $|x-1|$ représente la distance entre $x$ et 1.
    \end{itemize}
\end{itemize}

\subsection*{\color{myred} 2) Propriétés permettant de résoudre les équations et les inéquations du $1^{\text{ère}}$ degré}

Pour tout réel $a \geq 0$ et tout réel $x$ :

\begin{itemize}
    \item $|x| = a \iff x = a \text{ ou } x = -a$.
    \item $|x| \leq a \iff -a \leq x \leq a$.
    \item $|x| \geq a \iff x \leq -a \text{ ou } x \geq a$.
    \item Plus généralement : $|x - c| \leq a \iff c - a \leq x \leq c + a$ (cercle de centre $c$ et de rayon $a$ sur la droite réelle).
\end{itemize}

\begin{itemize}
    \item \textbf{Exemples :}
    \begin{itemize}
        \item Résoudre $|2x-1| = 3$ :
        \[
        2x-1 = 3 \quad \text{ou} \quad 2x-1 = -3
        \]
        Donc $2x = 4 \implies x=2$ ou $2x = -2 \implies x=-1$. Les solutions sont $S = \{-1, 2\}$.
        \item Résoudre $|x+2| \leq 1$ :
        \[
        -1 \leq x+2 \leq 1 \implies -3 \leq x \leq -1.
        \]
        Donc $S = [-3, -1]$.
        \item Résoudre $|3x-6| > 2$ :
        \[
        3x-6 < -2 \quad \text{ou} \quad 3x-6 > 2
        \]
        \[
        3x < 4 \implies x < \frac{4}{3} \quad \text{ou} \quad 3x > 8 \implies x > \frac{8}{3}.
        \]
        Donc $S = \left]-\infty, \dfrac{4}{3}\right[ \cup \left]\dfrac{8}{3}, +\infty\right[$.
    \end{itemize}
\end{itemize}

\vspace{0.5em}
\noindent \textbf{Remarque :} La valeur absolue est très utile pour exprimer des distances, des tolérances (erreurs), et intervient dans de nombreux domaines des mathématiques.

\end{document}